\scanpage{188}
% manuscript position: top
\mspage{104}
% manuscript position: 1–7
aux petites $\varinjlim$, il reste donc à voir qu'il commute aux $\varinjlim$ de cardinal $\leq\pi$ ; en vertu de 4.4.14.1. (marge page 98). Pour l'établir, on utilise le th. 4.4.13. b), en représentant un foncteur $\varphi:J\longrightarrow\operatorname{Ind}_\pi(C)$ \inserted{(où $\operatorname{card}\operatorname{Fl} J\leq\pi$)} sous forme indicielle par
% manuscript position: 8
\[
\Phi:I\times J\longrightarrow C,\tag*{(4.4.16.6)}
\]
% manuscript position: 9–10
$I$ étant un ens. ordonné grand devant $\pi$. On a donc
% manuscript position: 11
\[
\varphi(j)=\varinjlim_{i\in I}{}_{\operatorname{Ind}_\pi C}\Phi(i,j),
\]
% manuscript position: 12–13
et comme $g=\psi(f)$ commute aux $\varinjlim$ de type $I$ par ce qu'on a vu, on aura
% manuscript position: 14
\[
\deleted{g(\varinjlim{}_{\operatorname{Ind}_\pi C}\varphi(j))}
\]
% manuscript position: 15
\[
g(\varphi(j))=\varinjlim_{i\in I}{}_{(\unsure{\mathcal E|E})}g(\Phi(i,j))
\]
% manuscript position: 16
d'où
% manuscript position: 17
\[
\varinjlim_{j\in J}{}_{(\mathcal E)}g(\varphi(j))=\varinjlim_{j\in J}\varinjlim_{i\in I}
\underbrace{g(\Phi(i,j))}_{\substack{\Vert\\ f(\Phi(i,j))}},\tag*{(\(\ast\))}
\]
% manuscript position: right/17
\inserted{(limites \unsure{dans|ds} $\mathcal E$)}

% manuscript position: 18
D'autre part on a
% manuscript position: 19–20
\[
\begin{aligned}
\varinjlim_{j\in J}{}_{\operatorname{Ind}_\pi}\varphi(j)
 &\mathrel{\unsure{=|\simeq}}\varinjlim_{\unsure{J|j}\in J}\varinjlim_{i\in I}\Phi(i,j)\\
 &\simeq\varinjlim_{i\in I}\underbrace{\varinjlim_{j\in J}\Phi(i,j)}_{\substack{\text{\unsure{peut être|peut-être} calculé}\\\text{dans }C,\text{ soit}\\\Psi(i)\in\operatorname{Ob} C}},
\end{aligned}
\]
% manuscript position: 21
donc
% manuscript position: 22
\[
\varinjlim_{j\in J}{}_{\operatorname{Ind}_\pi C}\varphi(j)\mathrel{\unsure{=|\simeq}}\varinjlim_{i\in I}\unsure{\Psi|\Psi'}(i),\tag*{(\(\ast\ast\))}
\]
% manuscript position: foot
% archive sheet number: 187
